\cleardoublepage % memastikan bab baru dimulai di halaman ganjil (kanan)
\chapter{PENDAHULUAN}
\label{chap:pendahuluan}
% --- Latar Belakang ---
\section{Latar Belakang}
\label{sec:latar_belakang}

% ELEMEN 1: Kondisi/situasi + permasalahan
Sistem blockchain modern, seperti \textit{Ethereum}, telah bertransisi menggunakan mekanisme konsensus \textit{Proof-of-Stake} (PoS) yang keamanannya secara fundamental bertumpu pada sistem insentif ekonomi. Dalam mekanisme ini, validator yang bertindak jujur dan mengikuti aturan protokol akan diberikan imbalan (\textit{reward}), sedangkan validator yang terbukti melakukan kecurangan akan dikenakan penalti berupa pemotongan aset yang mereka pertaruhkan (\textit{slashing}) \parencite{buterin2020combining}. Permasalahan utama yang mengemuka saat ini adalah parameter pertahanan insentif tersebut bersifat statis dan hanya disetel secara manual di awal, sementara pihak penyerang terus beradaptasi. Dalam praktiknya, berbagai ancaman dari validator strategis ini telah terwujud secara nyata, mulai dari serangan \textit{selfish mining} \parencite{eyal2018majority}, serangan \textit{reorg} dan \textit{balancing}, (\textit{transaction censorship}), hingga eksploitasi \textit{Maximal Extractable Value} (MEV) untuk meraup keuntungan dari selisih harga transaksi pengguna lain \parencite{buterin2020combining}.

% ELEMEN 2: Urgensi / dampak negatif jika tidak diselesaikan
Persoalan parameter pertahanan statis yang berhadapan dengan penyerang yang terus berinovasi ini memiliki urgensi tinggi untuk segera ditangani, karena berdampak langsung terhadap aspek \textit{safety} dan \textit{liveness} dari jaringan \textit{blockchain}. Jika arsitektur pertahanan tidak pernah dievaluasi dan disesuaikan, terdapat dua risiko kritis yang dapat terjadi. Pertama, transaksi \textit{user} yang seharusnya telah mencapai tahap finalisasi dapat dibatalkan melalui serangan manipulasi rantai. Kedua, jaringan dapat mengalami kebuntuan konsensus sehingga transaksi baru tertahan tanpa kepastian. Mengingat total nilai aset yang dipertaruhkan di dalam jaringan \textit{Ethereum} telah mencapai angka puluhan miliar dolar, kedua risiko tersebut bukan sekadar potensi teoretis, melainkan ancaman nyata yang dapat mengurangi kepercayaan publik terhadap keandalan \textit{blockchain} sebagai infrastruktur keuangan yang aman.

% ELEMEN 3: Solusi yang telah/dapat diterapkan
Sejumlah pendekatan telah diusulkan dalam literatur untuk menganalisis dan memitigasi serangan pada mekanisme insentif \textit{blockchain}. Pendekatan awal berfokus pada analisis manual terhadap perilaku strategis penyerang, yang membuktikan bahwa kecurangan tetap dapat menguntungkan meskipun penyerang tidak menguasai mayoritas daya komputasi jaringan \parencite{eyal2018majority}. Perkembangan berikutnya memanfaatkan kecerdasan artifisial melalui riset SquirRL untuk mengotomasi proses ini; penggunaan \textit{Deep Reinforcement Learning} (DRL) terbukti lebih efektif dalam menemukan strategi serangan otomatis dibandingkan analisis manual \parencite{hou2021squirrl}. Di ranah yang lebih luas, kerangka kerja \textit{Multi-Agent Reinforcement Learning} (MARL) telah dieksplorasi secara matang untuk memodelkan interaksi dinamis antar-agen kompetitif \parencite{zhang2021multi}, serta telah diterapkan secara luas pada domain keamanan siber \parencite{nguyen2021deep}.berbagai kasus keamanan siber secara umum \parencite{nguyen2021deep}.

% ELEMEN 4: Kelemahan solusi -> dasar rumusan masalah TA
Meskipun berbagai upaya telah dilakukan, solusi-solusi yang ada masih meninggalkan celah penelitian yang mendasar. Analisis ancaman secara manual, seperti pada \textcite{eyal2018majority} tidak memiliki skalabilitas untuk menemukan pola varian serangan baru, sehingga prosesnya harus diulang dari awal setiap menghadapi ancaman berbeda. Di sisi lain, metode penemuan serangan otomatis, seperti SquirRL \parencite{hou2021squirrl} berfokus pada lingkungan (\textit{single-agent}) yang hanya melatih sisi penyerang, sementara parameter pertahanan protokol tetap dibiarkan statis selama proses evaluasi. Akibatnya, suatu desain parameter yang hanya efektif terhadap strategi serangan spesifik pada saat itu, dan rentan terhadap penyerang yang terus berevolusi. Lebih jauh, meskipun algoritma \textit{Multi Agent Reinforcement Learning} (MARL) telah matang \parencite{zhang2021multi}, pemanfaatannya sebagai pengendali yang adaptif untuk mengamankan konsensus PoS masih belum dieksplorasi dalam literatur. Kesenjangan inilah yang menjadi dasar pemikiran utama dalam merumuskan masalah tugas akhir ini, yaitu perlunya membangun arena penyerang-pertahanan pada konsensus PoS di mana agen pertahanan dapat turut belajar dan mengatur ulang respons protokol, seperti \textit{slashing} dan \textit{reward}, secara adaptif menggunakan pendekatan \textit{Multi-Agent Reinforcement Learning} (MARL).

\section{Rumusan Masalah}
\label{sec:rumusan_masalah}
Pokok persoalan yang menjadi fokus pada tugas akhir ini adalah evaluasi ketahanan protokol \textit{Proof-of-Stake} (PoS) terhadap serangan otomatis (sebagaimana dirintis oleh riset seperti SquirRL) yang selama ini dilakukan dengan membiarkan parameter pertahanan tetap statis sepanjang proses pelatihan agen penyerang. Akibatnya, kesimpulan bahwa suatu parameter pertahanan sudah cukup "aman" berpotensi menyesatkan, sebab kesimpulan tersebut hanya berlaku secara spesifik terhadap strategi serangan yang kebetulan ditemukan agen pada saat itu. Begitu agen penyerang diberi kesempatan untuk menyesuaikan strateginya lebih lanjut, misalnya melalui pelatihan tambahan atau variasi kondisi jaringan, tidak ada jaminan bahwa parameter pertahanan yang sama masih akan efektif dalam memitigasi serangan.

Urgensi penyelesaian masalah ini terletak pada konsekuensi finansial dan reputasional yang dapat timbul jika celah adaptasi penyerang ini luput terdeteksi sebelum protokol berjalan di jaringan produksi. Mengingat total nilai aset yang dipertaruhkan pada jaringan PoS berskala besar, seperti \textit{Ethereum}, berjumlah sangat signifikan, kesalahan estimasi dalam menilai ketahanan protokol terhadap serangan yang terus beradaptasi bukan sekadar persoalan akademis. Hal ini berpotensi menimbulkan kerugian ekonomi yang nyata bagi \textit{validator} jujur maupun pengguna jaringan, serta mengancam aspek keamanan (\textit{safety}) dan ketersediaan (\textit{liveness}) sistem secara keseluruhan.

Sebagai usulan solusi, tugas akhir ini memodelkan interaksi penyerang dan pertahanan pada konsensus PoS sebagai permainan multi-agen yang dilatih secara bersamaan melalui skema \textit{self-play} menggunakan \textit{Multi-Agent Reinforcement Learning} (MARL). Melalui pendekatan ini, agen penyerang dan agen pertahanan sama-sama diberikan ruang untuk beradaptasi sebelum ketahanan dari kebijakan protokol tersebut dinilai secara final. Alur pergeseran cara pandang ini, dari evaluasi parsial terhadap parameter statis menuju evaluasi komprehensif terhadap entitas penyerang yang adaptif.

Berdasarkan kerangka pemikiran tersebut, rumusan masalah dalam tugas akhir ini dijabarkan melalui tiga pertanyaan penelitian (\textit{Research Questions} / RQ) berikut:
\begin{itemize}
    \item \textbf{RQ1:} Sejauh mana formulasi \textit{Markov game}, yang mencakup ruang \textit{state}, \textit{action}, dan \textit{reward}, terbukti valid dan memadai untuk merepresentasikan dinamika adaptasi antara penyerang dan pertahanan pada konsensus PoS secara realistis tanpa melanggar spesifikasi dasar protokol?
    \item \textbf{RQ2:} Apakah kebijakan pertahanan yang dihasilkan melalui pelatihan MARL dengan skema \textit{self-play} menunjukkan tingkat ketahanan (\textit{robustness}) yang lebih baik terhadap agen penyerang adaptif, dibandingkan dengan penggunaan parameter pertahanan statis?
    \item \textbf{RQ3:} Bagaimana karakteristik strategi serangan yang berhasil ditemukan secara otomatis oleh agen \textit{Reinforcement Learning} di dalam arena simulasi, dan seberapa efektif kebijakan pertahanan adaptif dalam menekan profit relatif penyerang serta menekan tingkat pelanggaran \textit{safety} maupun \textit{liveness}?
\end{itemize}

\section{Tujuan}
\label{sec:tujuan}
Bertolak dari rumusan masalah di atas, tugas akhir ini bertujuan untuk:

\begin{enumerate}
	\item Merancang dan mengimplementasikan simulator konsensus \textit{Proof-of-Stake} yang mengacu pada protokol Gasper, dikalibrasi terhadap data \textit{on-chain Ethereum}, serta merumuskan interaksi antara penyerang dan pertahanan dalam arena tersebut sebagai \textit{Markov Game}(\textit{state, action}, dan \textit{reward}).
	
	\item Melatih agen pertahanan menggunakan algoritma \textit{Multi Agent Reinforcement Learning} (MARL), yaitu MAPPO dengan skema \textit{self-play} dan \textit{population-based training} agar mampu menyesuaikan parameter \textit{reward}, \textit{slashing}, dan \textit{checkpoint} finalitas, kemudian membandingkan ketahanannya terhadap kebijakan pertahanan statis pada populasi agen penyerang adaptif yang sama.
	
	\item Menganalisis pola strategi yang muncul dari kedua sisi agen selama pelatihan, serta mengukur dampaknya terhadap tingkat pelanggaran \textit{safety} dan \textit{liveness}, keuntungan relatif penyerang, dan biaya pertahanan berupa \textit{reward} yang hilang dari validator jujur.
\end{enumerate}

Kriteria keberhasilan tugas akhir ini ada dua. Pertama, simulator \textit{Proof-of-Stake} (PoS) yang dibangun harus mampu mereproduksi perilaku nominal jaringan (kondisi tanpa serangan) yang selaras secara statistik dengan data \textit{on-chain} sungguhan, sebagai syarat sebelum simulator layak dipakai untuk eksperimen lanjutan. Kedua, kebijakan pertahanan hasil MARL harus menunjukkan penurunan yang terukur pada setidaknya satu dari dua indikator berikut dibandingkan pertahanan statis, yaitu profit relatif penyerang, atau tingkat pelanggaran \textit{safety}/\textit{liveness}, ketika diuji terhadap populasi penyerang yang tidak dilihat agen pertahanan selama proses pelatihan.

\section{Ruang Lingkup, Asumsi, dan Batasan Masalah}

Supaya pelaksanaan tugas akhir ini realistis diselesaikan dalam jangka waktu dua semester, ditetapkan ruang lingkup dan batasan berikut:

\begin{enumerate}
	\item Protokol konsensus yang dimodelkan dibatasi pada Gasper (GHOST + Casper FFG), yaitu protokol yang dipakai \textit{Ethereum} dan menjadi acuan simulator pada tugas akhir ini. Protokol PoS lain, seperti \textit{Tendermint} pada \textit{Cosmos} atau \textit{Ouroboros} pada \textit{Cardano}, tidak menjadi objek penelitian ini, sebab mekanisme \textit{fork-choice} dan finalitasnya berbeda secara mendasar dan butuh pemodelan sendiri.
	
	\item Jenis serangan yang dimodelkan dalam arena dibatasi pada tiga 
    kategori yang telah diuraikan pada \ref{sec:latar_belakang}, 
    yaitu perilaku validator strategis sejenis \textit{selfish mining}, 
    serangan reorganisasi rantai (\textit{reorg}/\textit{balancing attack}), 
    dan \textit{transaction censorship} terhadap \textit{attestation} atau 
    proposal. Serangan jangka panjang (\textit{long-range attack}) dan 
    eksploitasi MEV sengaja tidak dimasukkan, karena keduanya membutuhkan 
    asumsi model dan skala simulasi yang berbeda. Kedua topik ini diusulkan 
    sebagai arah pengembangan lanjutan.
	
	\item Ruang aksi agen pertahanan dibatasi pada tiga parameter insentif protokol sebagaimana dirumuskan pada \ref{sec:tujuan}, yaitu laju \textit{reward per epoch}, tingkat keparahan atau ambang \textit{slashing}, dan parameter finalitas \textit{checkpoint}. Perubahan pada struktur algoritmik protokol itu sendiri, misalnya mengganti aturan \textit{fork-choice}, berada di luar ruang aksi agen pertahanan.
	
	\item Simulator yang dibangun adalah model ringan pada \textit{level slot} dan \textit{epoch}, bukan implementasi node \textit{Ethereum} yang lengkap, dan tidak dijalankan pada jaringan \textit{testnet} maupun \textit{mainnet} sungguhan. Asumsi ini diambil demi kelayakan komputasi, mengingat sumber daya yang tersedia terbatas pada CPU \textit{multi-core} dan GPU kelas menengah.
	
	\item Skala jaringan yang disimulasikan dibatasi pada kisaran ratusan sampai seribu \textit{validator}, bukan skala \textit{mainnet Ethereum} yang melibatkan ratusan ribu \textit{validator} aktif.
	
	\item Kalibrasi parameter simulator mengacu pada data \textit{on-chain} publik \textit{Ethereum}, diambil melalui dataset publik \textit{BigQuery} dan \textit{beaconcha.in}, serta spesifikasi resmi konsensus \textit{Ethereum} yang dipublikasikan di repositori \textit{consensus-specs}, sejalan dengan kriteria keberhasilan validasi simulator pada \ref{sec:tujuan}. Diasumsikan data historis tersebut cukup mewakili pola perilaku jaringan secara umum, meski tidak menutup kemungkinan adanya penyimpangan pada periode-periode tertentu.
\end{enumerate}

Ringkasan cakupan penelitian terhadap komponen utama arena penyerang dan pertahanan disajikan pada Tabel \ref{tab:scope_summary}.

\begin{table}[htbp]
\centering
\caption{Ruang lingkup penelitian pada setiap komponen arena \textit{Multi Agent Reinforcement Learning} (MARL)}
\label{tab:scope_summary}
\begin{tabular}{|p{3cm}|p{5.5cm}|p{5.5cm}|}
\hline
\textbf{Komponen} & \textbf{Termasuk Ruang Lingkup} & \textbf{Di Luar Ruang Lingkup} \\
\hline
Protokol & Gasper (GHOST + Casper FFG) & Tendermint, Ouroboros, protokol PoS lain \\
\hline
Jenis serangan & Selfish mining, reorg/balancing attack, censorship & Long-range attack, eksploitasi MEV \\
\hline
Aksi pertahanan & Reward rate, slashing, checkpoint finalitas & Perubahan aturan fork-choice \\
\hline
Simulator & Model ringan level slot/epoch & Full node, testnet, mainnet \\
\hline
Skala jaringan & Ratusan sampai ribuan validator & Skala mainnet (ratusan ribu validator) \\
\hline
Kalibrasi data & BigQuery, beaconcha.in, consensus-specs & Data proprietary atau non-publik \\
\hline
\end{tabular}
\end{table}

\section{Hipotesis Tugas Akhir}
Premis yang mendasari hipotesis tugas akhir ini disusun dari dua temuan pada literatur yang sudah dibahas pada \ref{sec:latar_belakang}. Pertama, Hou dkk. (2021) melalui SquirRL menunjukkan bahwa agen RL mampu menemukan strategi serangan yang lebih menguntungkan dibanding strategi rancangan manusia ketika dihadapkan pada parameter pertahanan yang tetap, sebuah kondisi yang menjadi pokok persoalan pada \ref{sec:rumusan_masalah}. Kedua, Zhang, Yang, dan Basar (2021) menunjukkan bahwa skema \textit{self-play} pada MARL memungkinkan suatu kebijakan menjadi tangguh terhadap lawan yang terus berubah, karena kebijakan tersebut dilatih menghadapi berbagai versi dirinya sendiri yang semakin kuat dari waktu ke waktu. Bila kedua temuan ini digabungkan, kebijakan pertahanan yang ikut dilatih untuk beradaptasi, sebagaimana diusulkan pada \ref{sec:tujuan}, semestinya menghasilkan ketahanan yang lebih baik dibanding pertahanan statis, ketika sama-sama dihadapkan pada penyerang adaptif. Dari premis tersebut, dirumuskan pasangan hipotesis berikut:

\begin{itemize}
	\item $H_0$: Tidak terdapat perbedaan yang signifikan pada metrik keamanan (tingkat pelanggaran \textit{safety}/\textit{liveness} dan profit relatif penyerang) antara kebijakan pertahanan hasil pelatihan MARL dengan \textit{self-play} dan kebijakan pertahanan dengan parameter statis, ketika keduanya diuji terhadap populasi penyerang adaptif yang sama.
	
	\item $H_1$: Kebijakan pertahanan hasil pelatihan MARL dengan \textit{self-play} menghasilkan tingkat pelanggaran \textit{safety}/\textit{liveness} yang lebih rendah dan menekan keuntungan relatif penyerang secara lebih signifikan dibandingkan kebijakan pertahanan dengan parameter statis, ketika diuji terhadap populasi penyerang adaptif yang sama. Namun, penurunan ini disertai kenaikan biaya pertahanan, berupa berkurangnya \textit{reward} yang diterima \textit{validator} jujur, yang perlu ditelaah lebih lanjut sebagai bentuk \textit{trade-off}, sejalan dengan metrik biaya pertahanan yang telah ditetapkan pada \ref{sec:tujuan}.
\end{itemize}

Keterkaitan antara premis literatur, rumusan masalah, dan pasangan hipotesis di atas dirangkum pada Tabel \ref{tab:hypothesis_basis}.

\begin{table}[htbp]
\centering
\caption{Keterkaitan premis literatur dengan rumusan hipotesis tugas akhir}
\label{tab:hypothesis_basis}
\begin{tabular}{|p{3.2cm}|p{5.8cm}|p{4.8cm}|}
\hline
\textbf{Premis} & \textbf{Temuan Literatur} & \textbf{Kaitan dengan Hipotesis} \\
\hline
Premis 1 & Hou dkk. (2021): agen RL menemukan strategi serangan lebih optimal terhadap pertahanan statis (SquirRL) & Mendasari $H_0$ sebagai kondisi \textit{default} yang diuji \\
\hline
Premis 2 & Zhang dkk. (2021): \textit{self-play} pada MARL menghasilkan kebijakan tangguh terhadap lawan yang terus beradaptasi & Mendasari $H_1$ sebagai dugaan arah perbaikan ketahanan \\
\hline
Gabungan & Pertahanan adaptif vs. pertahanan statis, diuji pada populasi penyerang adaptif yang sama & Menjawab RQ2 dan RQ3 pada Rumusan Masalah \\
\hline
\end{tabular}
\end{table}


\section{Metodologi}

Tugas akhir ini menempuh pendekatan eksperimen simulasi kuantitatif, dengan tahapan kerja sebagai berikut:

\begin{enumerate}
	\item \textbf{Investigasi dan perumusan masalah.} Tahap ini mencakup 
    penelusuran kondisi keamanan konsensus PoS saat ini sebagaimana 
    diuraikan pada \ref{sec:latar_belakang}, dilanjutkan dengan 
    \textit{systematic literature review} (SLR) menggunakan \textit{query} 
    berikut pada basis data Scopus: 
    \texttt{TITLE-ABS-KEY(("proof of stake" OR blockchain OR consensus) 
    AND "reinforcement learning" AND (attack OR defense OR security OR 
    "selfish mining"))}. Hasil SLR, termasuk peta klaster literatur dan 
    posisi kontribusi tugas akhir ini terhadap celah yang diidentifikasi 
    pada Rumusan Masalah, dituangkan secara komprehensif pada Bab II.
	
	\item \textbf{Perancangan simulator dan formulasi permainan.} Membangun simulator konsensus PoS pada level slot dan \textit{epoch} (mencakup pemilihan komite, aturan \textit{fork-choice}, perhitungan \textit{reward}, dan mekanisme \textit{slashing}) sesuai batasan protokol Gasper pada Ruang Lingkup, dikalibrasi terhadap data \textit{on-chain} publik, lalu merumuskan \textit{state}, \textit{action}, dan \textit{reward space} untuk kedua agen sebagai \textit{Markov game} guna menjawab RQ1.
	
	\item \textbf{Validasi simulator.} Membandingkan keluaran simulator pada kondisi nominal, yaitu tanpa keberadaan agen penyerang, terhadap statistik agregat jaringan \textit{Ethereum} sungguhan. Simulator baru dianggap layak dipakai untuk tahap berikutnya jika selisihnya berada dalam rentang yang dapat diterima, sesuai kriteria keberhasilan pertama pada \ref{sec:tujuan}.
	
	\item \textbf{Reproduksi baseline serangan.} Melatih agen penyerang RL \textit{single-agent}, mengikuti pendekatan \textit{SquirRL}, pada simulator yang sudah dibangun, melawan pertahanan dengan parameter statis. Tahap ini berfungsi sebagai titik pembanding terhadap kondisi pertahanan statis yang menjadi pokok persoalan pada \ref{sec:rumusan_masalah}, sebelum arena multi-agen yang lebih kompleks dibangun pada semester berikutnya.
	
	\item \textbf{Pembangunan arena MARL.} Mengembangkan lingkungan pelatihan multi-agen menggunakan kerangka kerja terbuka (PettingZoo dan RLlib), lalu melatih kedua agen secara bergantian menggunakan algoritme MAPPO dengan skema \textit{self-play}, dilanjutkan dengan \textit{population-based training} yang melibatkan beberapa varian kebijakan penyerang sekaligus, sebagaimana ditetapkan pada \ref{sec:tujuan} poin kedua.
	
	\item \textbf{Eksperimen dan evaluasi.} Menguji kebijakan pertahanan hasil MARL berhadapan dengan kebijakan pertahanan statis, keduanya dihadapkan pada populasi penyerang adaptif yang sama, untuk menguji hipotesis $H_0$ dan $H_1$. Pengujian diulang dengan beberapa \textit{seed} acak agar hasil dapat diuji signifikansinya secara statistik. Metrik yang digunakan meliputi tingkat pelanggaran \textit{safety} dan \textit{liveness}, profit relatif penyerang, biaya pertahanan, serta skor \textit{robustness} lintas populasi penyerang.
	
	\item \textbf{Analisis dan penarikan kesimpulan.} Menganalisis pola strategi yang muncul dari kedua sisi agen, mengaitkan temuan tersebut dengan rumusan masalah dan hipotesis yang diajukan, serta menyusun kesimpulan dan saran pengembangan lanjutan, termasuk kemungkinan perluasan ke serangan jangka panjang dan eksploitasi MEV yang tidak tercakup pada batasan penelitian ini.
\end{enumerate}

Keterkaitan setiap tahapan metodologi dengan rumusan masalah yang dijawab dirangkum pada Tabel \ref{tab:methodology_rq}.

\begin{table}[htbp]
\centering
\caption{Keterkaitan tahapan metodologi dengan bagian \ref{sec:rumusan_masalah} tugas akhir}
\label{tab:methodology_rq}
\begin{tabular}{|p{0.6cm}|p{5cm}|p{2cm}|p{5cm}|}
\hline
\textbf{No.} & \textbf{Tahapan} & \textbf{Menjawab} & \textbf{Keluaran Utama} \\
\hline
1 & Investigasi dan perumusan masalah & Latar Belakang & Peta literatur, posisi kontribusi (Bab II) \\
\hline
2 & Perancangan simulator dan formulasi permainan & RQ1 & Simulator Gasper, formulasi Markov game \\
\hline
3 & Validasi simulator & Kriteria 1 (Tujuan) & Simulator tervalidasi terhadap data \textit{on-chain} \\
\hline
4 & Reproduksi baseline serangan & Konteks $H_0$ & Baseline agen penyerang \textit{single-agent} \\
\hline
5 & Pembangunan arena MARL & RQ1, RQ2 & Kebijakan pertahanan hasil \textit{self-play} \\
\hline
6 & Eksperimen dan evaluasi & RQ2, RQ3, $H_0$/$H_1$ & Metrik \textit{safety}, \textit{liveness}, profit, \textit{robustness} \\
\hline
7 & Analisis dan penarikan kesimpulan & Seluruh RQ & Simpulan dan saran pengembangan \\
\hline
\end{tabular}
\end{table}

\section{Sistematika Penulisan}
\label{sec:sistematika}

Laporan tugas akhir ini disusun ke dalam beberapa bab dengan sistematika penulisan sebagai berikut:

\begin{enumerate}
    \item \textbf{Bab I Pendahuluan:} Berisi gambaran umum tugas akhir, meliputi latar belakang, rumusan masalah, tujuan penelitian, ruang lingkup dan batasan masalah, hipotesis penelitian, metodologi yang digunakan, serta sistematika penulisan laporan.
    \item \textbf{Bab II Studi Literatur:} Berisi kajian pustaka dan dasar teori yang relevan dengan topik tugas akhir, termasuk tinjauan terhadap serangan strategis pada konsensus \textit{blockchain}, penerapan \textit{reinforcement learning}, kerangka kerja \textit{Multi-Agent Reinforcement Learning} (MARL), dan identifikasi kesenjangan (\textit{gap}) penelitian.
    \item \textbf{Bab III Analisis Masalah:} Berisi analisis dan evaluasi ketahanan sistem pertahanan PoS statis saat ini, formulasi \textit{Markov game} untuk interaksi agen penyerang-pertahanan, serta rancangan skema pendekatan MARL \textit{self-play} sebagai solusi yang diusulkan.
    \item \textbf{Bab IV Desain Konsep Solusi:} Berisi penjelasan terperinci mengenai rancangan arsitektur dan parameter simulator konsensus PoS, rancangan ruang agen (pendefinisian \textit{state}, \textit{action}, dan \textit{reward}), serta alur skema pelatihan \textit{population-based training}.
    \item \textbf{Bab V Rencana Selanjutnya:} Berisi penjabaran langkah-langkah implementasi, rancangan eksperimen dan skenario pengujian \textit{baseline} (statis) versus MARL, penetapan metrik evaluasi (\textit{safety, liveness, profit, robustness}), serta analisis risiko penelitian beserta mitigasinya.
\end{enumerate}